\clearpage
\section[Tujuan dan Manfaat]{}
\sectionopener{cgreen}{Kenapa Ini Penting}{Tujuan dan Manfaat}{Atlas dirancang untuk satu tujuan sederhana: satu platform, seluruh alur kerja tim, tanpa biaya berlapis.}

\subsection{Tujuan pengembangan}

\begin{enumerate}[leftmargin=1.4em]
  \item Menyatukan manajemen proyek (PMO) dan komunikasi tim (chat, suara) dalam satu aplikasi, sehingga tim tidak perlu berpindah-pindah antara alat yang terpisah.
  \item Menyediakan platform yang dapat di-\textit{self-host} sepenuhnya, tanpa batas jumlah pengguna, proyek, maupun fitur yang dikunci di balik paket berbayar.
  \item Mendukung otentikasi tingkat institusi (SSO SAML/OIDC, OAuth, OTP telepon, tautan ajaib) secara gratis, agar organisasi dengan sistem identitas sendiri tetap bisa mengintegrasikan Atlas tanpa biaya tambahan.
  \item Memberi admin kendali penuh atas konfigurasi instansi (Godmode) tanpa perlu menyentuh kode atau berkas konfigurasi server.
  \item Membuktikan kelayakan solusi melalui pemakaian produksi nyata, bukan sekadar purwarupa kompetisi.
\end{enumerate}

\subsection{Manfaat bagi penerima manfaat}

\noindent
\begin{minipage}[t]{0.485\linewidth}
\featurecard{cgreen}{Laboratorium dan himpunan kampus}{Mengelola puluhan proyek dan puluhan anggota dalam satu ruang kerja, tanpa anggaran langganan SaaS yang biasanya mustahil dipenuhi organisasi mahasiswa nirlaba.}
\end{minipage}\hfill
\begin{minipage}[t]{0.485\linewidth}
\featurecard{cgreen}{Startup dan tim kecil tahap awal}{Mendapatkan perangkat PMO dan komunikasi setara produk berbayar sejak hari pertama, sehingga anggaran terbatas bisa dialihkan ke pengembangan produk inti.}
\end{minipage}

\vspace{0.5em}
\noindent
\begin{minipage}[t]{0.485\linewidth}
\featurecard{cgreen}{Komunitas dan organisasi nirlaba}{Menjalankan koordinasi program dan relawan dengan kendali penuh atas data anggota, tanpa bergantung pada kebijakan atau harga sepihak vendor asing.}
\end{minipage}\hfill
\begin{minipage}[t]{0.485\linewidth}
\featurecard{cgreen}{Ekosistem digital Indonesia}{Menambah alternatif perangkat lunak kolaborasi bersumber terbuka yang dapat dipelajari, dimodifikasi, dan dikontribusikan balik oleh komunitas developer lokal.}
\end{minipage}

\vspace{1.1em}
\begin{quotebreakbox}
{\bricolagesb\fontsize{16pt}{21pt}\selectfont\color{white} Dampak yang paling meyakinkan bukan proyeksi, melainkan yang sudah terjadi.}\\[0.5em]
{\fontsize{10pt}{14pt}\selectfont\color{white!70!black} Sejak dipakai penuh oleh Laboratorium MGM, seluruh koordinasi proyek internal (perencanaan tugas, linimasa, dan komunikasi harian 65+ anggota) berjalan di Atlas, menggantikan kombinasi spreadsheet dan obrolan tersebar yang dipakai sebelumnya, tanpa satu rupiah pun biaya langganan bulanan.}
\end{quotebreakbox}

\clearpage
